% Generated draft: requires mathematical and source review.
\documentclass[11pt]{article}
\usepackage[margin=25mm]{geometry}
\usepackage{fontspec}
\setmainfont{Latin Modern Roman}
\usepackage{amsmath,amssymb,mathtools,unicode-math}
\setmathfont{Latin Modern Math}
\usepackage{enumitem,needspace,xurl,hyperref}
\setlength{\parindent}{0pt}
\setlength{\parskip}{5pt}
\setlength{\emergencystretch}{3em}
\newcommand{\N}{\mathbb{N}}
\newcommand{\Z}{\mathbb{Z}}
\newcommand{\Q}{\mathbb{Q}}
\newcommand{\R}{\mathbb{R}}
\newcommand{\C}{\mathbb{C}}
\newcommand{\F}{\mathbb{F}}
\newcommand{\A}{\mathbb{A}}
\newcommand{\PP}{\mathbb{P}}
\newcommand{\E}{\mathbb{E}}
\newcommand{\eps}{\varepsilon}
\newcommand{\dd}{\,\mathrm{d}}
\begin{document}
\section*{UnsolvedMath ID: 701702}
\textbf{Source problem code:} catalogue entry
\subsection{Definitions and mathematical statement}
The supplied statement refers to Problem 17.57 for the definition of \(\mathrm{CT}(\Z)\), but that definition is not included. Consequently, the ambient group and its action cannot be specified completely from this record. The statement presupposes an action of \(\mathrm{CT}(\Z)\) on \(\Z\).

Relative to that intended action, for a subgroup \(H\leq \mathrm{CT}(\Z)\) define
\[
H\cdot n=\{h\cdot n:h\in H\},\qquad
\mathcal O_H=\{H\cdot n:n\in\Z\}.
\]
The subgroup \(H\) is finitely generated if \(H=\langle F\rangle\) for some finite subset \(F\subseteq H\). A set of orbit representatives is a subset \(R\subseteq\Z\) satisfying
\[
\forall O\in\mathcal O_H,\qquad |R\cap O|=1.
\]
The source asks whether every finitely generated subgroup \(H\leq\mathrm{CT}(\Z)\) satisfies
\[
|\mathcal O_H|<\infty
\quad\text{or}\quad
\exists R\subseteq\Z\;\Bigl[
(\forall O\in\mathcal O_H,\ |R\cap O|=1)
\ \land\ R\text{ has positive density}\Bigr].
\]
Here ``positive density'' retains the source wording: the record does not specify the density convention on \(\Z\). This is therefore a conditional transcription of the question, pending the definition in Problem 17.57 and clarification of density, rather than a fully specified conjecture.
\subsection{Short English statement}
Must every finitely generated subgroup of \(\mathrm{CT}(\Z)\) either have finitely many orbits on \(\Z\) or admit a set containing exactly one representative from each orbit and having positive density? The supplied record omits the definition of \(\mathrm{CT}(\Z)\) and the intended density convention.
\subsection{Sources}
\begin{itemize}
\item Supplied local UnsolvedMath registry record: permanent ID \(701702\), exact source problem code \texttt{catalogue entry}, title ``Kourovka Notebook Problem 20.44.'' The registry labels the problem open; this is supplied metadata, not an independently verified status assessment.
\item Attributed source: \emph{The Kourovka Notebook (Unsolved Problems in Group Theory)}, Problem 20.44, proposed by S. Kohl in 2022. The supplied background names the 20th issue (2022) and the 21st issue edition (2026), with identifier \texttt{arXiv:1401.0300v48}. The supplied collection description names E. I. Khukhro and V. D. Mazurov as editors.
\item The statement explicitly refers to Problem 17.57 for the definition of \(\mathrm{CT}(\Z)\). That referenced definition is absent from the supplied record. The cited source and cross-reference have not been independently consulted or verified.
\end{itemize}
\end{document}
